\documentclass[10pt,a4paper]{amsart}
\usepackage[margin=19mm]{geometry}
\usepackage{amsmath,amssymb,mathtools}
\usepackage[scr=rsfs,cal=euler]{mathalfa}
\usepackage{xfrac}
\usepackage[unicode,hidelinks,bookmarks=false]{hyperref}
\newcommand{\ie}{i.e.\ }
\newcommand{\cf}{cf.\ }
\newcommand{\resp}{respectively\ }
\newcommand{\Subsection}{\subsection}
\providecommand{\nobreakdash}{\nobreak\hbox{-}}
\setlength{\emergencystretch}{2em}
\multlinegap0pt
\usepackage{bm}
\usepackage{bbm}
\usepackage{dsfont}
\usepackage{xspace}
\usepackage{cancel}
\usepackage{mathtools}
\usepackage{amsthm}
\makeatletter
\@ifundefined{theorem}{\newtheorem{theorem}{Theorem}}{}
\@ifundefined{lemma}{\newtheorem{lemma}[theorem]{Lemma}}{}
\@ifundefined{remark}{\newtheorem{remark}[theorem]{Remark}}{}
\makeatother
\makeatletter
\newcommand{\E}{\mathbb{E}} % Expectation E
\DeclareMathOperator{\N}{\mathbb{N}} % Natural numbers
\newcommand{\be}{\begin{equation}}
\expandafter\ifx\csname breakablealgorithm\endcsname\relax
\newenvironment{breakablealgorithm}
  {% \begin{breakablealgorithm}
   \begin{center}
     \refstepcounter{algorithm}% New algorithm
     \hrule height.8pt depth0pt \kern2pt% \@fs@pre for \@fs@ruled
     \renewcommand{\caption}[2][\relax]{% Make a new \caption
       {\raggedright\textbf{\ALG@name~\thealgorithm} ##2\par}%
       \ifx\relax##1\relax % #1 is \relax
         \addcontentsline{loa}{algorithm}{\protect\numberline{\thealgorithm}##2}%
       \else % #1 is not \relax
         \addcontentsline{loa}{algorithm}{\protect\numberline{\thealgorithm}##1}%
       \fi
       \kern2pt\hrule\kern2pt
     }
  }{% \end{breakablealgorithm}
     \kern2pt\hrule\relax% \@fs@post for \@fs@ruled
   \end{center}
  }
\fi
\newcommand{\ee}{\end{equation}}
\newcommand{\p}{\mathbb{P}} % Probability P
\makeatother
\title[On the largest common subtree of uniform attachment trees]{On the largest common subtree of uniform attachment trees}
\author{Johannes B\"aumler, C\'eline Kerriou, Bas Lodewijks, James Martin, Emil Powierski, Mikl\'os Z. R\'acz, Anirudh Sridhar}
\date{Source: arXiv:2609.30098v1 (CC BY 4.0)}
\begin{document}
\maketitle

\noindent\emph{Source and licence.} On the largest common subtree of uniform attachment trees. Version arXiv:2609.30098v1 (CC BY 4.0), distributed under the Creative Commons Attribution 4.0 International licence, as recorded in the arXiv licence metadata for this exact version and shown on the abstract page https://arxiv.org/abs/2609.30098v1. Full rights notice: licenses/arxiv-2609.30098-CC-BY-4.0.txt.\par\smallskip

\noindent\textit{Editorial scope.} Counted unit: the Lemma whose source label is \texttt{lemma:pdembed} in the article (source0.tex lines 3444--3449, source label \texttt{lemma:pdembed}), delivered together with the article's own complete original proof of it (source0.tex lines 3451--3540, one \texttt{proof} environment of 90 lines). No mathematics is written, repaired or re-derived: statement and proof are the article's own text line for line. Required context retained uncounted: remark:pois diri (lines 3384--3389). Every definition, display and label of those lines is retained unchanged. Omitted from this package and not redistributed: 1--3383, 3390--3443, 3450--3450, 3541--4633. Nothing inside a retained excerpt is omitted or altered: every retained body is delivered exactly as the article's lines are written, so no omission receipt of any kind accompanies this package. Citations: the retained excerpts carry no citation command at all, so no bibliography entry of the article is transcribed and no reference is redistributed. Reference adaptation: none was needed and none was made; every reference inside the delivered unit resolves inside it, so no unresolved reference or citation is produced. Printed slips kept literal: no printed word, symbol, hypothesis or number of the counted statement or of its proof was altered. Layout and class: the article's own class is \texttt{article} with packages including amsmath, amssymb, amsbsy, amsthm, dsfont, graphicx, xcolor, tikz, caption, subcaption, enumerate, ulem, cases, geometry; the wrapper is the lane's pinned \texttt{amsart} editorial wrapper at 10pt on A4, which restates the theorem-like environments the retained text needs and adds the article's own macro definitions that the retained text needs (\texttt{\textbackslash E}, \texttt{\textbackslash N}, \texttt{\textbackslash be}, \texttt{\textbackslash ee}, \texttt{\textbackslash p}), with extra wrapper packages (bm, bbm, dsfont, xspace, cancel, mathtools). The delivered numbering is the unit's own. Assets: no retained span contains a figure environment, a graphics-inclusion command, a PDF-page inclusion, a table or a TikZ picture; no image file is copied into this package, the package declares no asset dependency and no image file is redistributed with it. Licence: version arXiv:2609.30098v1 of this article is distributed under the Creative Commons Attribution 4.0 International licence, as recorded in the arXiv licence metadata for this exact version and shown on the abstract page https://arxiv.org/abs/2609.30098v1. A full rights notice is delivered at licenses/arxiv-2609.30098-CC-BY-4.0.txt. Characters: every retained excerpt is pure ASCII, so the delivered engine, XeLaTeX with its default Latin Modern text font, reports no missing character.
\begin{remark}\label{remark:pois diri}
It is readily observed that, for $v\in \N^N$ and $d\in [N]$ such that $v_1+\cdots+v_d=N$ and $v_{d+1}=\ldots=v_N=0$, 
\begin{equation*}
    \p\left(w(N)=v \right) = \prod_{i=1}^{d} \frac{1}{\sum_{j=i}^{d} v_j}.
\end{equation*}
\end{remark}

\begin{lemma}\label{lemma:pdembed}
    Let $K\geq 1$ be an integer and let $v\in \N^K$ with $|v|_1 \geq 2$. Then,
    \be 
    \E[|\{v\rightarrow w(|v|_1)\}|]\leq 0.9^{K}.
    \ee 
\end{lemma}

\begin{proof} 
We start with the case $K=1$. In this case, we have that $|\{v\rightarrow w(|v|_1)\}| \in \{0,1\}$, with $|\{v\rightarrow w(|v|_1)\}| = 1$ if and only if $w(|v|_1) = (|v|_1)$, which occurs with probability $\frac{1}{|v|_1} \leq \frac{1}{2} \leq 0.9^1$.

From here on, we will assume that $K \geq 2$.
Suppose that, for each $i\in \N$, the vector $v$ contains exactly $z_i\in\N_0$ entries equal to $i$ such that $\sum_{i=1}^\infty z_i = K$ and let $z = (z_1, z_2, z_3, \dots)$.
The statement is clear when $z_1=K$ and $z_i = 0$ for $i\geq 2$, since in this case
\begin{equation*}
    \E[|\{v\rightarrow w(|v|_1)\}|] = \p \left( w(|v|_1) = v  \right) = \frac{1}{K!} \leq 0.9^K
\end{equation*}
for all $K \geq 2$. Thus, we will assume that $z_1 < K$ for the rest of the proof.
There are 
\be 
\binom{K}{z} = \frac{K!}{\prod_{i=1}^{\infty} z_i!}
\ee 
many ways to reorder the vector $v$. We write $\hat{v}$ for the reordering of $v$ in decreasing order, that is, such that $\hat{v}_1\geq \hat{v}_2\geq \ldots  \geq \hat{v}_K.$ For each reordering $\overline{v}$ of $v$ we have that 
\begin{align*}
	\p\left(w\left(|v|_1\right) = \overline{v} \right)
	\leq
	\p\left(w\left(|v|_1\right) =\hat{v} \right),
\end{align*}
as for any vector $v=(v_1,\ldots,v_K)$, by Remark \ref{remark:pois diri},
\begin{align}\label{eq:equality decreasing}
    \p\left(w\left(|v|_1\right) = v \right) = \prod_{i=1}^{K} \frac{1}{\sum_{j=i}^{K} v_j},
\end{align}
which is maximal when the values $(v_1,\ldots, v_K)$ are in decreasing order. Using a union bound over all $\binom{K}{z}$ many reorderings of $v$, we directly get that 
\be 
\p\left(w\left(|v|_1\right) \sim v \right) \leq \binom{K}{z} \p\left(w\left(|v|_1\right) = \hat{v} \right).
\ee
Thus,
\begin{align}\label{eq:isomorphis expect}
	 \notag \E \left[| \{ v \rightarrow w \left(|v|_1\right) \}|\right]  &= \p \left( v \sim w \left(|v|_1\right) \right)
	 \E \left[| \{ v \rightarrow w \left(|v|_1\right) \}| \ \Big| \ v \sim w \left(|v|_1\right) \right]\\
	 & \notag
	 \leq
	\binom{K}{z}
	\p\left(w\left(|v|_1\right) =  \hat{v}  \right)
	\prod_{i=2}^{\infty} z_i!\\
	&
	\leq 
	\frac{K!}{z_1! } \p\left(w\left(|v|_1\right) =  \hat{v}  \right).
\end{align}
From \eqref{eq:equality decreasing} we also get that
\begin{align*}
    \p\left(w\left(|v|_1\right) =  \hat{v}  \right) 
    &= 
    \prod_{i=1}^{K} \frac{1}{\sum_{j=i}^{K}  \hat{v}_j} 
    = 
    \frac{1}{z_1!} \prod_{i=1}^{K-z_1} \frac{1}{\sum_{j=i}^{K}  \hat{v}_j}\\
    &
    \leq
    \frac{1}{z_1!} \prod_{i=1}^{K-z_1} \frac{1}{z_1 + 2(K-z_1 - i+ 1)}
    =
    \frac{1}{z_1!} \prod_{\ell=1}^{K-z_1} \frac{1}{z_1 + 2\ell},
\end{align*}
where we used variable transform $\ell=K-z_1-i+1$ for the last equality. Thus, we obtain for all $K\geq 2$ that 
\begin{align}\label{eq:outdegree bound}
	\E \left[| \{ v \rightarrow w \left(|v|_1\right) \}|\right]\leq \frac{K!}{z_1! } \p\left(w\left(|v|_1\right) =  \hat{v} \right)
    \leq \frac{K!}{z_1!} \frac{1}{z_1!} \prod_{\ell=1}^{K-z_1} \frac{1}{z_1 + 2\ell}
	= 
	\frac{1}{z_1!} \prod_{\ell=1}^{K- z_1} \frac{z_1+\ell}{z_1+2\ell}.
\end{align}
It thus remains to prove that we can bound the right-hand side from above by $0.9^K$. To this end, we define $s:=K-z_1>0$. In the following, we consider three different cases for combinations of $s$ and $z_1$, and we prove the desired upper bound for all three cases.

    \noindent
    \textbf{Case 1:}  $s\geq 2z_1$. Then, for $\ell \geq \lceil s/2 \rceil$ one has $\frac{z_1+\ell}{z_1+2\ell} \leq \frac{2}{3}$. Furthermore, the condition $s\geq 2z_1$ implies that
    $s = 2s/3 + s/3 \geq 2s/3 + 2z_{1}/3 = 2K/3$. 
    % \begin{align*}
    %     s = \frac{2}{3}s + \frac{1}{3}s \geq \frac{2}{3}s + \frac{2}{3}z_1 = \frac{2}{3} K .
    % \end{align*}
    Combining these two observations, we get that
    \begin{equation*}
        \prod_{\ell=1}^{s}\frac{z_1+\ell}{z_1+2\ell} \leq \prod_{\ell=\lceil s/2 \rceil}^{s}\frac{z_1+\ell}{z_1+2\ell} \leq \left(\frac{2}{3}\right)^{s/2} \leq \left(\frac{2}{3}\right)^{K/3} \leq 0.9^K.
    \end{equation*}
    \textbf{Case 2:}   $z_1> 1, s\leq 2z_1$. 
    It follows that $z_1 \geq K/3$, since 
    $3 z_{1} = 2 z_{1} + z_{1} \geq s + z_{1} = K$. 
    % \begin{align*}
    %     3 z_1 = 2z_1 + z_1 \geq s+z_1 = K.
    % \end{align*}
    % We thus get that
    Thus, 
    \begin{align*}
        \frac{1}{z_1!}\leq \frac{1}{2^{z_1-1}} \leq \frac{1}{2^{z_1/2}} \leq \frac{1}{2^{K/6}} \leq 0.9^K.
    \end{align*}
\textbf{Case 3:} $z_1=1, s=1$. Here we directly have that
\begin{align*}
    \frac{1}{z_1!} \prod_{\ell=1}^{s} \frac{z_1+\ell}{z_1+2\ell} = \frac{2}{3} \leq 0.9^2 = 0.9^K,
\end{align*}
which concludes the proof of Lemma~\ref{lemma:pdembed}.
\end{proof}

\end{document}
